\chapter{分论概说}

% Source PDF physical page 215; see audit/source_coverage.json.
\begin{EditionOpeningPanel}
分论复习的方法论：第一，抓大放小。“大”包括：财产犯罪，人身犯罪，贪污贿赂犯罪，其他章节的重点罪名。第二，掌握罪名的行为公式。第三，善于对比总结。第四，记住司法解释的特殊规定，同时务必掌握正常原理。
\end{EditionOpeningPanel}

\section{罪名与罪状}

[问题]甲欲拐卖男童，误将15周岁的乙女当作男童，予以拐卖。对甲如何处理？（拐卖妇女案，2021年试题）

\subsection*{一、罪名}

罪名有类罪名和具体罪名之分。例如，分则第三章第五节“金融诈骗罪”就是类罪名，其下的保险诈骗罪就是具体罪名。

1.单一罪名。特征：一个条文只表述一个罪名，该罪名不能分拆使用。例如，第232条只规定了故意杀人罪，这个罪名不能分拆使用。

2.选择性罪名。特征：条文规定的犯罪行为和犯罪对象多种多样。例如，第280条规定的伪造、变造、买卖国家机关公文、证件、印章罪。

（1）类型：第一，只规定了多个选择行为。例如，“引诱、容留、介绍卖淫罪”。第二，只规定了多个选择对象。例如，“拐卖妇女、儿童罪”。第三，同时规定了多个选择行为和选择对象。例如，“非法制造、买卖、运输、邮寄、储存枪支、弹药、爆炸物罪”。

（2）特点：既可以概括使用，也可以分拆使用。

[注意1]\CJKunderwave{行为对象不是一个对象时，也概括使用，不数罪并罚}。

例1，就拐卖妇女、儿童罪而言，当只拐卖妇女时，定拐卖妇女罪；当只拐卖儿童时，定拐卖儿童罪；当既拐卖妇女，又拐卖儿童时，也只定拐卖妇女、儿童罪，不以拐卖妇女罪和拐卖儿童罪并罚。

例2，王某引诱甲女卖淫，容留乙女卖淫，介绍丙女卖淫（3女都已不是幼女），只定引诱、容留、介绍卖淫罪。

% Source PDF physical page 216; see audit/source_coverage.json.
例3，王某出售假的人民币，购买假的美元，运输假的欧元，只定出售、购买、运输假币罪。

[注意2］对选择性对象之间的对象认识错误，按客观定。例如，本节开头“拐卖妇女案”中，按照客观定，甲构成拐卖妇女罪既遂。

3.概括罪名。特征：条文规定的行为类型多种多样。

（1）罪名只能概括使用，不能分拆使用。例如，信用卡诈骗罪包括了使用伪造的信用卡、使用作废的信用卡、冒用他人信用卡、恶意透支四种行为类型。行为人只实施了恶意透支行为，应定信用卡诈骗罪，不能分拆成“恶意透支罪”。

（2）同时实施了多个行为类型，只定一个罪名，不数罪并罚。例如，行为人将信用卡诈骗罪的四种类型全部实施了，也只定一个信用卡诈骗罪。

\subsection*{二、罪状}

1.简单罪状。例如，刑法第232条规定“故意杀人的”，构成故意杀人罪。\CJKunderwave{注意}：简单罪状不违反罪刑法定原则的明确性要求，因为没有必要详细描述。

2.叙明罪状。这是指详细描述行为类型的罪状。

3.引证罪状。例如，刑法第115条第1款规定了放火罪，第2款规定“过失犯前款罪的，处…”，该款引用第1款罪状，用来表达失火罪。这便是引证罪状。

4.空白罪状。例如，刑法第345条第2款规定“违反森林法的规定，滥伐森林或者其他林木，数量较大的”，构成滥伐林木罪。该规定没有详细描述滥伐林木罪的具体行为类型，做了留白，需要参照森林法的规定。这便是空白罪状。

[总结］空白罪状不等于引证罪状。引证罪状，引用的是刑法自身条文。空白罪状，引用的是行政法等条文。（2022年试题）

\section{注意规定与法律拟制}

[问题］狗蛋带领小兄弟与狗剩一帮人聚众斗殴。过程中，狗蛋过失致人死亡。对狗蛋如何处理？（聚众斗殴案）

注意规定和法律拟制是解读刑法分则条文的一把钥匙，务必掌握。\footnote{关于注意规定与法律拟制的详细内容，参见张明楷：《刑法分则的解释原理》（第二版）（下），中国人民大学出版社2011年版，第622\textasciitilde{}682页。}

\subsection*{一、注意规定}

注意规定，是指在刑法已作基本规定的前提下，提示司法人员注意、以免司法人员忽略的规定。其特点是：并不改变基本规定的内容，只是对基本规定内容的重申；即使不设置该规定，遇到此类情形也应按照基本规定处理。例如，司法解释规定，携带挪用的公款潜逃的，按照贪污罪论处。该规定就是注意规定。因为携带挪用的公款潜逃的，由于具有了非法占有目的，本身就构成贪污罪。这条规定的出台只是提醒法官不要将这种情形定为挪用公款罪。即使没有这条规定，遇到携带挪用的公款潜逃的情形，也应按照贪污罪论处。

% Source PDF physical page 217; see audit/source_coverage.json.
因此，判断一项法律条文是不是注意规定的方法是：假设取消这项条文，遇到条文中的事项，处理结论是否仍应如此。如果是，表明该条文只具有提示意义，那么该条文就是注意规定。例如，第253条第2款规定，邮政工作人员私自开拆、隐匿、毁弃邮件、电报并从中窃取财物的，按照盗窃罪论处。假设没有该款规定，遇到这类情况，仍应按照盗窃罪处罚，所以该款只具有提示作用，属于注意规定。

[总结]刑法分则中常见的注意规定：

1.利用计算机盗窃。

第287条规定：“利用计算机实施金融诈骗、盗窃、贪污、挪用公款、窃取国家秘密或者其他犯罪的，依照本法有关规定定罪处罚。”此条属于注意规定。即使没有该条规定，这类行为也构成具体实施的犯罪，如金融诈骗罪犯罪、盗窃罪等。

2.第382条规定的是贪污罪，第3款规定：“与前两款所列人员勾结，伙同贪污的，以共犯论处。”该款是注意规定，因为即使没有该款规定，对此情形也应以贪污罪共犯论处。正因如此，第385条规定的受贿罪虽没有类似规定，但是与国家工作人员勾结，伙同受贿的，也应以受贿罪共犯论处。

3.第310条规定的是窝藏、包庇罪，第2款规定：“犯前款罪，事前通谋的，以共同犯罪论处。”例如，甲与乙事前通谋共同盗窃，由甲去偷，事后由乙负责窝藏包庇甲。乙构成盗窃罪的共同犯罪，不构成窝藏、包庇罪。这表明，该款是注意规定。正因如此，第312条的掩饰、隐瞒犯罪所得罪虽没有类似规定，但是与犯罪人事前通谋的，也应以事前通谋之罪的共同犯罪论处。例如，甲与乙事前通谋共同盗窃，由甲去偷，偷来的财物，由乙负责销售。乙构成盗窃罪的共同犯罪，不构成掩饰、隐瞒犯罪所得罪。

\subsection*{二、法律拟制}

法律拟制，是将原本不符合某种规定的行为也按照该规定处理。例如，刑法第269条规定，犯盗窃、诈骗、抢夺罪，为窝藏赃物、抗拒抓捕或者毁灭罪证而当场使用暴力或者以暴力相威胁的，依照刑法第263条（抢劫罪）的规定定罪处罚。该条将盗窃、诈骗、抢夺的情形按照抢劫罪处理。如果没有该条规定，只能定盗窃罪（或诈骗罪、抢夺罪）与故意伤害罪，数罪并罚。

因此，判断一项法律条文是不是法律拟制的方法是，\CJKunderwave{假设取消该条文，遇到条文中的}\CJKunderwave{事项，处理结论是否仍会如此。如果不是，表明该条文是特殊规定，不具有普遍适用性，}\CJKunderwave{那么该条文就是法律拟制}。例如，第289条规定，聚众“打砸抢”，毁坏或者抢走公私财物的，对首要分子，依照抢劫罪论处。假设取消该条规定，遇到聚众“打砸抢”，毁坏财物的，应定故意毁坏财物罪，而不应定抢劫罪。这表明该条规定将不属于抢劫罪的情形拟制成抢劫罪，因此属于法律拟制。简言之，\CJKunderwave{注意规定属于重复强调，无特殊之处；法律拟}\CJKunderwave{制属于特殊规定，不能推而广之}。

[总结]刑法分则中常考的法律拟制：

1.非法拘禁罪\ensuremath{\rightarrow}故意伤害罪、故意杀人罪

第238条第2款后一句规定，非法拘禁过程中，使用暴力致人伤残、死亡的，按照故意伤害罪、故意杀人罪论处。这里的“使用暴力”是指非法拘禁的实行行为之外的暴力行为。这里的“致人伤残、死亡”，是指过失所为。也即，该规定将过失致人重伤、死亡的

% Source PDF physical page 218; see audit/source_coverage.json.
行为拟制为故意伤害罪、故意杀人罪。

2.刑讯逼供罪、暴力取证罪\ensuremath{\rightarrow}故意伤害罪、故意杀人罪

第247条规定，刑讯逼供、暴力取证，致人伤残、死亡的，依照故意伤害罪、故意杀人罪论处。这里的致人伤残、死亡，是指过失所为。也即，该规定将过失致人重伤、死亡的行为拟制为故意伤害罪、故意杀人罪。（2010年第58题）

3.虐待被监管人罪一故意伤害罪、故意杀人罪

第248条规定，虐待被监管人，致人伤残、死亡的，依照故意伤害罪、故意杀人罪论处。该规定的原理与上述第247条完全相同。（2012年第16题）

4.聚众斗殴罪一故意伤害罪、故意杀人罪

第292条第2款规定，聚众斗殴，致人重伤、死亡的，依照故意伤害罪、故意杀人罪论处。这里的“致人重伤、死亡”，是指过失所为。也即，该规定将过失致人重伤、死亡的行为拟制为故意伤害罪、故意杀人罪。例如，本节开头的“聚众斗殴案”中，对狗蛋以故意杀人罪论处。

5.聚众“打砸抢”\ensuremath{\rightarrow}故意伤害罪、故意杀人罪、抢劫罪

第289条规定，聚众“打砸抢”，致人伤残、死亡的，依照故意伤害罪、故意杀人罪论处；毁坏或者抢走公私财物的，除判令退赔外，对首要分子，依照抢劫罪论处。首先，其中致人伤残、死亡既包括故意所为，也包括过失所为。就故意所为而言，该条就是注意规定；就过失所为而言，该条就是法律拟制。其次，将毁坏财物按照抢劫罪论处，属于法律拟制；将抢夺走财物定抢劫罪，属于法律拟制；将抢劫走财物定抢劫罪，属于注意规定。

6.抢夺罪一抢劫罪

第267条第2款规定，携带凶器抢夺的，依照抢劫罪论处。这属于法律拟制，将一种特殊的抢夺行为拟制为抢劫罪。

7.盗窃、诈骗、抢夺罪一抢劫罪

第269条规定，犯盗窃、诈骗、抢夺罪，为窝藏赃物、抗拒抓捕或者毁灭罪证而当场使用暴力或者以暴力相威胁的，依照本法第263条（抢劫罪）的规定定罪处罚。该条属于法律拟制，将盗窃罪（诈骗罪、抢夺罪）+故意伤害罪（致人轻伤的前提下而言）拟制为抢劫罪。

8.信用卡诈骗罪一盗窃罪

第196条第3款规定，盗窃信用卡并使用的，依照盗窃罪论处。因为盗窃信用卡如果不使用，就作无罪处理，所以处罚的重点是使用行为。如果在机器上使用，构成盗窃罪。此时该条就属于注意规定，因为在机器上使用信用卡本身就是盗窃。如果对自然人使用（如对银行服务人员使用），构成信用卡诈骗罪。此时该条就属于法律拟制，将信用卡诈骗罪拟制成盗窃罪。

9.为卖淫、通风报信一窝藏、包庇罪

第362条规定，旅馆业、饮食服务业、文化娱乐业、出租汽车业等单位的人员，在公安机关查处卖淫、活动时，为违法犯罪分子通风报信，情节严重的，依照第310条（窝藏、包庇罪）的规定定罪处罚。如果行为人通风报信的是卖淫、等违法行为，则该条是法律拟制。因为卖淫、者是违法分子，不是犯罪分子，而窝藏、包庇罪要求的

% Source PDF physical page 219; see audit/source_coverage.json.
行为对象是犯罪分子。如果行为人通风报信的是强迫卖淫等犯罪行为，则该条是注意规定。

[统计]（1）法律拟制来的故意杀人罪，共有5个来源罪名：非法拘禁罪，刑讯逼供罪，暴力取证罪，虐待被监管人罪，聚众斗殴罪。（2）法律拟制来的抢劫罪，共有3个：第267条第2款（携带凶器抢夺定抢劫），第269条（事后转化抢劫），第289条（聚众“打砸抢”中故意毁坏财物、抢夺财物，定抢劫）。

\OriginalPanel{典型真题}

关于《刑法》分则条文的理解，下列哪些选项是错误的？（2011年·卷二·58题\footnote{[答案]AB。}

A.即使没有《刑法》第二百六十九条的规定，对于犯盗窃罪，为毁灭罪证而当场使用暴力的行为，也要认定为抢劫罪

B.即使没有《刑法》第二百六十七条第二款的规定，对于携带凶器抢夺的行为也应认定为抢劫罪

C.即使没有《刑法》第一百九十六条第三款的规定，对于盗窃信用卡并在ATM取款的行为，也能认定为盗窃罪

D.即使没有《刑法》第一百九十八条第四款的规定，对于保险事故的鉴定人故意提供虚假的证明文件为他人实施保险诈骗提供条件的，也应当认定为保险诈骗罪的共犯

